\subsection*{MT-M37. Compatible nonnegative indentations cannot carry
the external prime shells}

The normalized replacements in MT-M36 do not by themselves
establish an orbital decomposition.  Here is a stronger
obstruction to one particular repair: a \emph{single}
nonnegative indentation normalized both globally and on
the initial local subgroup must vanish outside that subgroup.
This is exactly the simultaneous normalization offered in
Masullo's Proposition 12.2.2.4, original pp.~512--513.
An exact defect identity quantifies the obstruction
when the local normalization is relaxed.

\paragraph{The two normalizations force support.}
Let $G$ be an abelian group, $H\subset G$ a countable
subgroup, and $G_0\subset G$ a subgroup such that
\[
                 G=HG_0,\qquad H_0=H\cap G_0.
 \tag{M37.1}
\]
Suppose $b:G\to\mathbf R_{\ge0}$ satisfies
\[
 \sum_{h\in H}b(hx)=1\quad(x\in G),\qquad
 \sum_{h\in H_0}b(hx)=1\quad(x\in G_0).
 \tag{M37.2}
\]
Then
\[
                       b(y)=0\qquad(y\notin G_0).
 \tag{M37.3}
\]
For its averaged product
\[
                     B_b(x,y)=\sum_{h\in H}b(hx)b(hy),
 \tag{M37.4}
\]
one has
\[
                  B_b(x,y)=0\qquad(x^{-1}y\notin G_0).
 \tag{M37.5}
\]
In particular, if $G_0\subset G_k\subset G_{k+1}$ and
$\alpha_k=\mathbf1_{G_k}$, then
\[
       (\alpha_{k+1}(y)-\alpha_k(y))B_b(x,y)=0
                                      \quad(x\in G_0,\ y\in G).
 \tag{M37.6}
\]

\begin{proof}
For $x\in G_0$, subtract the second identity in (M37.2)
from the first.  Nonnegativity gives
\[
                       \sum_{h\in H\setminus H_0}b(hx)=0,
 \tag{M37.7}
\]
so every term is zero.  Any $y\notin G_0$ has the form
$hx$ with $h\in H$ and $x\in G_0$ by (M37.1).
That $h$ cannot belong to $H_0$, so (M37.7) proves
(M37.3).  If a term in (M37.4) is nonzero, both
$hx$ and $hy$ belong to $G_0$; hence $x^{-1}y\in G_0$.
This proves (M37.5).  The series (M37.4) is well-defined:
each factor is at most one by (M37.2), so its nonnegative
sum is at most $\sum_hb(hx)=1$.
Finally, for $x\in G_0$, either $y\notin G_0$ and
(M37.5) applies, or $y\in G_0\subset G_k$ and the
indicator difference is zero.  This proves (M37.6).
\end{proof}

The use of $hx$ instead of $h^{-1}x$ makes no difference
to the normalization, since inversion permutes each
subgroup.  The nonnegativity assumption is indispensable
to the termwise implication in (M37.7); it is part of
the source's definition of an indentation.

\paragraph{Exact defect and external orbital mass.}
Keep (M37.1) and the global normalization in (M37.2),
but do not impose the local normalization.  For $x\in G_0$
and a coset $c\in H/H_0$, define
\[
 A_c(x)=\sum_{h\in c}b(hx),\qquad
 \delta(x)=1-A_{H_0}(x).
 \tag{M37.8}
\]
Then $A_c\ge0$, $\sum_cA_c=1$, and $0\le\delta\le1$.
The total mass assigned to external translates is
\[
 \begin{split}
 W_{\mathrm{ext}}(x)
   &:=\sum_{\eta\in H\setminus H_0}B_b(x,\eta x)\\
   &=1-\sum_{c\in H/H_0}A_c(x)^2.
 \end{split}
 \tag{M37.9}
\]
In particular
\[
       2\delta(x)(1-\delta(x))
          \le W_{\mathrm{ext}}(x)
          \le2\delta(x)-\delta(x)^2\le2\delta(x).
 \tag{M37.10}
\]

\begin{proof}
Tonelli for nonnegative countable sums gives
\[
 \sum_{\eta\in H}B_b(x,\eta x)
    =\sum_{h\in H}\sum_{\eta\in H}b(hx)b(h\eta x)=1.
 \tag{M37.11}
\]
If $\eta$ ranges only over $H_0$, then $h\eta$ runs
over the same coset of $H_0$ as $h$.  Summing coset
by coset gives
\[
                   \sum_{\eta\in H_0}B_b(x,\eta x)
                              =\sum_c A_c(x)^2.
 \tag{M37.12}
\]
Subtracting proves (M37.9).  The identity coset has
mass $1-\delta$ and the other cosets have total mass
$\delta$, so
\[
       0\le\sum_{c\ne H_0}A_c^2\le\delta^2.
 \tag{M37.13}
\]
Substitution into (M37.9) gives both bounds in (M37.10).
\end{proof}

Here is the precise consequence for varying indentations.
Suppose $G$ carries a Haar measure, $b$ is measurable,
$D\subset G_0$ is a
measurable fundamental domain for $H$, and
$d\in L^1(D)$ is a prescribed quotient diagonal.
For any measurable selection of external translates,
its integral after multiplication by $B_b$ has absolute
value at most
\[
                    2\int_D\delta(x)|d(x)|\,d\mu_G(x).
 \tag{M37.14}
\]
Indeed the selected nonnegative weights are bounded by
$W_{\mathrm{ext}}$, and the quotient diagonal is the
same at $x$ and $\eta x$.  Equation (M37.10) proves
the bound; absolute convergence follows from
$W_{\mathrm{ext}}\le1$ and $d\in L^1$.
Thus if a sequence $b_j$ has defects $\delta_j\to0$
almost everywhere on $D$, every such external integral
tends to zero for the fixed diagonal, by dominated
convergence.  If the diagonals also vary, the same
conclusion follows either from
\[
       \|\delta_j\|_\infty\to0,\quad
                       \sup_j\|d_j\|_1<\infty,
 \tag{M37.15}
\]
or from $d_j\to d$ in $L^1$ and
$\delta_j\to0$ almost everywhere.  In the latter case
\[
 \int_D\delta_j|d_j|
       \le\|d_j-d\|_1+\int_D\delta_j|d|\longrightarrow0.
 \tag{M37.16}
\]
These are explicit sufficient conditions, not an
assertion that a varying source kernel meets them.

\paragraph{A trace-norm criterion for the varying diagonals.}
The second alternative has a useful operator formulation.
Fix a Hilbert realization $\mathcal H=L^2(D,\mu)$.
There is a canonical linear contraction
\[
 \mathcal D:\mathcal S_1(\mathcal H)\longrightarrow L^1(D,\mu),
 \qquad
 \operatorname{tr}(M_fT)=\int_D f\,\mathcal D(T)\,d\mu
                   \quad(f\in L^\infty(D,\mu)),
 \tag{M37.20}
\]
where $M_f$ denotes multiplication and $\mathcal S_1$ the
trace-class operators.  In particular
$\int_D\mathcal D(T)\,d\mu=\operatorname{tr}T$.
This is the trace diagonal; no pointwise restriction of an
arbitrary measurable two-variable kernel is being asserted.

\begin{proof}
Write a singular-value expansion
$T=\sum_n s_n\,u_n\otimes\overline{v_n}$, with
$\sum_n s_n=\|T\|_1$ and unit vectors $u_n,v_n$.
Set $\mathcal D(T)=\sum_n s_nu_n\overline{v_n}$.
Cauchy--Schwarz gives
$\|u_n\overline{v_n}\|_{L^1}\le1$, so the series converges
in $L^1$ with norm at most $\|T\|_1$.
For every bounded $f$, the rank-one trace formula and
trace-norm convergence give (M37.20).
Bounded measurable functions separate $L^1$ functions;
therefore (M37.20) makes the definition independent of
the expansion and proves linearity.  Applying the same
bound to $T_1-T_2$ proves contraction.
\end{proof}

Consequently, if $\|T_j-T\|_1\to0$ in this fixed realization,
$d_j=\mathcal D(T_j)$, and $\delta_j\to0$ almost everywhere,
every selected external orbital contribution in (M37.14)
tends to zero, with the explicit bound
\[
 2\|T_j-T\|_1+
       2\int_D\delta_j\,|\mathcal D(T)|\,d\mu
                         \longrightarrow0.
 \tag{M37.21}
\]
Indeed $\|d_j-\mathcal D(T)\|_{L^1}\le\|T_j-T\|_1$,
and the second term tends to zero by dominated convergence.
The assertion also applies to a fixed closed Hilbert subspace,
or to a closed subquotient, realized isometrically by $J$ in
this $L^2$ space: replace $T_j$ by $JT_jJ^*$, which preserves
trace norms and traces because it adds only zero singular
values.  If $w$ is measurable with $0<w<\infty$ almost
everywhere, a weighted $L^2(D,w\mu)$ realization can first
be conjugated by the unitary $Uu=\sqrt w\,u$ into $L^2(D,\mu)$.
The trace identity gives
$\mathcal D_\mu(UTU^{-1})=w\mathcal D_{w\mu}(T)$,
so (M37.21) uses the density for the fixed reference measure
after that conjugation.
This formulation requires that the proposed geometric diagonal
has actually been identified with the fixed trace diagonal.
Trace class alone does not identify it with a chosen version
of a measurable kernel on a null diagonal.

\paragraph{The source's compatible-indentation construction.}
Original Proposition 12.2.2.4, p.~512, assumes a second
countable locally compact abelian $G$, a discrete $H$
with counting measure, an increasing exhaustive family
of open subgroups $G_k$, and $H_k=H\cap G_k$.
Among its stated hypotheses are a fundamental domain
$D\subset G_0$ and the analogous fundamental-domain
property for each $G_k/H_k$.  In particular
$G=HD\subset HG_0$, which gives (M37.1).
Its asserted indentation $b$ is global and has
$\alpha_kb$ an $H_k$-indentation on $G_k$ for every
$k$.  At $k=0$ this is precisely the second
normalization in (M37.2).  The first is the global
indentation condition (12.2.1.2), p.~505.

Therefore every nonnegative indentation with the
asserted simultaneous normalization satisfies
(M37.3)--(M37.6).  This conclusion uses only the
two normalizations and $G=HG_0$; none of the additional
differentiation or domain-regularity hypotheses is
invoked.  The source's own construction on p.~513
explicitly extends an $H_0$-indentation on $G_0$
by zero to $G$, consistent with (M37.3).
In particular its tail integrand (12.2.2.26), evaluated
at $x\in D\subset G_0$ and $y=h^{-1}x$ outside $G_k$,
is already zero for each $k$.  The proposition's
zero-tail limit is not refuted by this observation.

\paragraph{The actual rational external-prime shells.}
Take
\[
 G=\mathbb I_{\mathbf Q},\quad H=\mathbf Q^\times,
 \quad G_0=\mathbf R^\times\times
                           \widehat{\mathbf Z}^\times,
 \quad H_0=\{1,-1\}.
 \tag{M37.17}
\]
Every idele can have its finite valuations canceled
by a rational number, so $G=HG_0$ directly.
Let $T_0=\varnothing$ and let $T_k$ exhaust all rational
primes, adding one prime at a time.  Define
\[
 G_k=\mathbf R^\times\times
         \prod_{p\in T_k}\mathbf Q_p^\times\times
         \prod_{p\notin T_k}\mathbf Z_p^\times.
 \tag{M37.18}
\]
These are open subgroups; adding one prime has quotient
$\mathbf Z$ by its valuation.  The standard section
$\tau(t)$ and every $g\in E_{\mathbf Q}^1=
\{1\}\times\widehat{\mathbf Z}^\times$ belong to $G_0$.
For \emph{any} single nonnegative $b$ satisfying the
global $\mathbf Q^\times$ normalization and the local
$\{1,-1\}$ normalization on $G_0$, (M37.6) gives
\[
 (\alpha_{k+1}-\alpha_k)(\xi g\tau(t))
                  B_b(\tau(t),\xi g\tau(t))=0
 \quad(\xi\in\mathbf Q^\times,\ g\in E_{\mathbf Q}^1,\ t>0).
 \tag{M37.19}
\]
This is pointwise, before integrating or taking any
kernel limit.  Multiplying by an additional radial
cutoff or by a pulled-back measurable kernel does
not change the vanishing wherever that kernel is
defined.  Thus no such single compatible indentation
can turn the source's external shell in (11.3.2.1)
into its nonzero local Poisson value, exhibited by
the explicit test (M36.7)--(M36.8).

The dependency distinction is precise.  The proof of
Theorem 11.3.2.1 on original p.~495 actually cites
Proposition 12.2.2.2 for its indented-kernel trace
interpretation, and on p.~496 claims an approximation
by the local orbital finite-rank kernels.  It does
\emph{not} there explicitly cite Proposition 12.2.2.4.
The latter is the later available construction of a
single indentation compatible with the whole local
filtration.  MT-M37 proves that substituting that
compatible construction for the invalid literal
section average of MT-M36 cannot supply the asserted
external-shell trace.  It does not invent a source
citation or refute the zero-tail statement itself.

\paragraph{What a varying-level repair must address.}
The family $b_t$ of MT-M36 drops the local normalization
at the initial subgroup and can redistribute mass to
the $p$ and $p^{-1}$ translates; its defect is exactly
$\delta=t$ on the standard domain, and it attains the
lower bound $2t(1-t)$ in (M37.10).  Thus there is no
obstruction here to all nonnegative global indentations.
The obstruction is to imposing the two normalizations
on the \emph{same} function.

A level-dependent family could instead normalize only
at its current subgroup and need not have zero defect
at the original $G_0$.  Such a replacement must specify
the resulting orbital weights and justify its kernel
and trace limits; (M37.14)--(M37.16) give concrete
tests for any claim that the initial defect disappears.
No boundedness, $L^1$ convergence, or invariant orbital
allocation is inferred from the existence of the
individual indentations.  The Hecke-GRH endpoint
$M_H$ remains conditional.

\paragraph{Audit scope.}
Original pp.~512--513 were freshly rendered and read
in full for the hypotheses and extension-by-zero
construction.  Original pp.~495, 505 and 508 were
visually checked in MT-M36, and the complete supplied
proof text on pp.~495--496 was read as a dependency
map.  The new support and defect statements above
are fully proved here and require no GRH input.
